\chapter{Volatility Models}\label{ch:qm:volatility-models}

A risk report measures each currency's volatility by the equally weighted standard deviation of its last 250 daily returns. On 11 November
2022 the euro rose 3.5\% against the dollar: it was the first ECB reference rate (set around 14:10 CET) after the US consumer price index for October, released
at 8:30 New York time on 10 November, showed annual inflation slowing to 7.7\%. At the close of 2 November 2023, 250 fixings later, that return leaves the window, and the
EUR/USD volatility the report carries into the next day falls from 0.538\% to 0.496\%, by 7.8\%, on a day when the euro rose 1.2\%. A model
that let old news fade gradually would have reported a rise. This chapter is about the models that do: the ARCH and GARCH family and their
exponentially weighted special case, \omterm{def:qm:volatility-models:rv}{realised variance} from intraday data and the heterogeneous autoregression that forecasts it, and the loss
functions and tests that decide which forecast is better when the truth is never observed.

\section{The conditional-heteroskedasticity family}

\begin{definition}[Volatility clustering, conditional heteroskedasticity]\label{def:qm:volatility-models:clustering}
\emph{Volatility clustering}\index{volatility clustering} is the tendency of large returns, of either sign, to follow large returns, and small
ones small: squared returns are positively autocorrelated although returns are nearly not. A return series has \emph{conditional
heteroskedasticity}\index{conditional heteroskedasticity} when its variance given the past, $h_t = \Var(r_t \mid \mathcal F_{t-1})$, varies over time.
\end{definition}

\begin{definition}[ARCH and GARCH models, volatility persistence]\label{def:qm:volatility-models:garch}
In an \emph{ARCH model}\index{ARCH model} (Engle, 1982), $r_t = \sqrt{h_t}z_t$ with iid standardised $z_t$ and $h_t = \omega + \sum_{i=1}^q\alpha_ir_{t-i}^2$.
The \emph{GARCH model}\index{GARCH model} (Bollerslev, 1986) adds lagged variances; GARCH(1,1) is $h_t = \omega + \alpha r_{t-1}^2 + \beta h_{t-1}$ with $\omega > 0$,
$\alpha, \beta \ge 0$. Its \emph{volatility persistence}\index{volatility persistence} is $\alpha + \beta$, the rate at which a variance shock decays.
\end{definition}

\begin{proposition}[Stationarity and half-life of GARCH(1,1)]\label{prop:qm:volatility-models:garch}
If $\alpha + \beta < 1$, the GARCH(1,1) return is weakly stationary with unconditional variance $\bar h = \omega/(1 - \alpha - \beta)$, and the forecast of the
variance $k$ days ahead is $\E_t[h_{t+k}] = \bar h + (\alpha + \beta)^{k-1}(h_{t+1} - \bar h)$. A shock to the variance halves in $\ln\frac12/\ln(\alpha + \beta)$
days.
\end{proposition}

\begin{proof}
$\E_t[r_{t+k}^2] = \E_t[h_{t+k}]$, so $\E_t[h_{t+k+1}] = \omega + (\alpha + \beta)\E_t[h_{t+k}]$, a linear recursion with fixed point $\bar h$ and rate $\alpha + \beta$.
Taking unconditional expectations gives $\E[r^2] = \bar h$ when $\alpha + \beta < 1$.
\end{proof}

The parameters are estimated by maximising the likelihood of the returns given the recursion, with normal or Student-$t$ innovations. When the
normal likelihood is used although the innovations are not normal, the \omterm{def:qm:estimation:estimator}{estimator} is a \omterm{def:qm:estimation:kl}{quasi-maximum likelihood} \omterm{def:qm:estimation:estimator}{estimator}: it remains consistent and
asymptotically normal under regularity conditions, and its \omterm{def:qm:estimation:estimator}{standard errors} must be the sandwich of chapter~11.

\begin{definition}[EWMA volatility]\label{def:qm:volatility-models:ewma}
The \emph{EWMA volatility}\index{EWMA volatility} with decay $\lambda$ is $h_t = \lambda h_{t-1} + (1 - \lambda)r_{t-1}^2$: an exponentially weighted average of past squared
returns. RiskMetrics (1996) set $\lambda = 0.94$ for daily data.
\end{definition}

EWMA is GARCH(1,1) with $\omega = 0$, $\alpha = 1 - \lambda$ and $\beta = \lambda$: integrated, $\alpha + \beta = 1$, so it has no unconditional variance to revert to and
its forecast at every horizon is today's value. With $\lambda = 0.94$ a return's weight halves in 11 days. The equally weighted window is the other extreme:
every return in the window has the same weight, and then none.

\begin{definition}[Leverage effect, GJR-GARCH model]\label{def:qm:volatility-models:gjr}
The \emph{leverage effect}\index{leverage effect} is the tendency of volatility to rise more after falls than after rises of the same size, strong for equity
indices. The \emph{GJR-GARCH model}\index{GJR-GARCH model} (Glosten, Jagannathan and Runkle, 1993) captures it with $h_t = \omega + (\alpha + \gamma\mathbf 1_{r_{t-1} <
0})r_{t-1}^2 + \beta h_{t-1}$; its persistence is $\alpha + \beta + \gamma/2$ for symmetric innovations.
\end{definition}

Fitted to 7\,098 daily EUR/USD returns (1999--2026, in percent), the three models give:
\begin{center}
\footnotesize
\setlength{\tabcolsep}{4pt}
\begin{tabular}{@{}lccccccc@{}}
\toprule
 & $\omega$ & $\alpha$ & $\beta$ & $\gamma$ & $\nu$ & $\alpha + \beta + \gamma/2$ & \omterm{def:qm:stochastic-differential-equations:ou}{half-life} \\
\midrule
GARCH, normal & 0.0010 & 0.029 (0.003) & 0.968 (0.004) & --- & --- & 0.9973 & 255 \\
GARCH, Student $t$ & 0.0007 & 0.030 (0.003) & 0.968 (0.003) & --- & 7.2 (0.6) & 0.9984 & 445 \\
GJR, Student $t$ & 0.0006 & 0.027 (0.005) & 0.969 (0.004) & 0.006 (0.005) & 7.2 (0.6) & 0.9986 & 482 \\
\bottomrule
\end{tabular}
\end{center}
with sandwich \omterm{def:qm:estimation:estimator}{standard errors} in parentheses. The Student-$t$ likelihood is higher by 141 log points for one extra parameter; the asymmetry is not
significant (a currency pair has no natural ``down'' direction, unlike an equity index); and the persistence is so close to one that a shock to EUR/USD
volatility takes well over a year to halve: the \omterm{def:qm:linear-time-series:longmem}{long memory} of volatility of chapter~17, seen through a short-memory model
(\cref{fig:qm:volatility-models:history}).

\begin{omfigure}
\begin{tikzpicture}
\begin{axis}[width=12cm,height=5cm,xlabel={year},ylabel={daily \% (absolute return, volatility)},
  tick label style={font=\small},label style={font=\small},xmin=1999,xmax=2026.8,ymin=0,ymax=3,grid=major,grid style={gray!25},
  xtick={2000,2005,2010,2015,2020,2025},xticklabel style={/pgf/number format/set thousands separator={}},
  legend style={font=\footnotesize,at={(0.98,0.96)},anchor=north east,fill=white,draw=none},table/col sep=comma]
\addplot[only marks,mark=*,mark size=0.3pt,black!35] table[x=year,y=absr]{figdata/methods/18-volatility-models/history.csv};
\addplot[omDef,very thick] table[x=year,y=garch]{figdata/methods/18-volatility-models/history.csv};
\legend{absolute daily return (every fifth day),GARCH-$t$ volatility}
\end{axis}
\end{tikzpicture}

\omcaption{Daily EUR/USD absolute returns and the conditional volatility of the fitted GARCH(1,1) with Student-$t$ innovations, 1999--2026. The
clusters (2000, 2008--2009, 2011, 2015, 2022) are what the model's persistence of 0.9984 describes. Data: ECB euro reference rates (source: ECB
statistics).}\label{fig:qm:volatility-models:history}
\end{omfigure}

The three \omterm{def:qm:estimation:estimator}{estimators} treat one large return very differently (\cref{fig:qm:volatility-models:impulse}). A 3.5\% day adds to the variance forecast an amount
that decays at 0.9984 a day under the fitted GARCH, at 0.94 a day under EWMA, and not at all under the window, which carries it undiminished for 250
days and then drops it at once. That drop is the hook's ghost (\cref{fig:qm:volatility-models:ghost}): on the day it happened, the window's volatility fell
7.8\%, while the GARCH volatility rose from 0.424\% to 0.465\% and the EWMA from 0.410\% to 0.490\%, both reacting to the day's 1.2\% move.

\begin{omfigure}
\begin{tikzpicture}
\begin{semilogyaxis}[width=10cm,height=5cm,xlabel={days after a 3.5\% return},ylabel={extra variance ($\%^2$)},
  tick label style={font=\small},label style={font=\small},xmin=0,xmax=300,ymin=1e-4,ymax=1,grid=major,grid style={gray!25},
  legend columns=3,legend style={font=\footnotesize,at={(0.5,-0.3)},anchor=north,draw=none,fill=none,
  /tikz/every even column/.append style={column sep=8pt}},table/col sep=comma]
\addplot[omDef,very thick] table[x=day,y=garch]{figdata/methods/18-volatility-models/impulse.csv};
\addplot[omThm,very thick,dashed] table[x=day,y=ewma]{figdata/methods/18-volatility-models/impulse.csv};
\addplot[black,very thick,dotted,const plot] table[x=day,y=window]{figdata/methods/18-volatility-models/impulse.csv};
\legend{{GARCH-$t$, fitted},{EWMA, 0.94},250-day window}
\end{semilogyaxis}
\end{tikzpicture}

\omcaption{How long one 3.5\% daily return stays in each variance forecast (the extra variance over a background of 0.5\% days). GARCH lets it decay slowly
(\omterm{def:qm:stochastic-differential-equations:ou}{half-life} 445 days), EWMA quickly (11 days), and the equally weighted window keeps it whole for 250 days and then drops it in one step. Data: the
chapter's tutorial, from the fitted parameters.}\label{fig:qm:volatility-models:impulse}
\end{omfigure}

\begin{omfigure}
\begin{tikzpicture}
\begin{axis}[width=11cm,height=5cm,xlabel={trading days relative to 2 November 2023},ylabel={daily volatility (\%)},
  tick label style={font=\small},label style={font=\small},xmin=-300,xmax=60,ymin=0.3,ymax=1.2,grid=major,grid style={gray!25},
  legend columns=3,legend style={font=\footnotesize,at={(0.5,-0.3)},anchor=north,draw=none,fill=none,
  /tikz/every even column/.append style={column sep=8pt}},table/col sep=comma]
\addplot[black,very thick,const plot] table[x=k,y=window]{figdata/methods/18-volatility-models/ghost.csv};
\addplot[omDef,very thick] table[x=k,y=garch]{figdata/methods/18-volatility-models/ghost.csv};
\addplot[omThm,thick,dashed] table[x=k,y=ewma]{figdata/methods/18-volatility-models/ghost.csv};
\legend{250-day window,GARCH-$t$,{EWMA, 0.94}}
\end{axis}
\end{tikzpicture}

\omcaption{EUR/USD volatility forecasts around the day the return of 11 November 2022 left the 250-day window (day 0, 2 November 2023). The window jumps up
by the same amount when the return enters (day $-250$) and falls by 7.8\% when it leaves; the GARCH and EWMA forecasts respond to the market, not to the
calendar. Data: ECB euro reference rates (source: ECB statistics).}\label{fig:qm:volatility-models:ghost}
\end{omfigure}

\section{Realised volatility}

\begin{definition}[Integrated variance, realised variance]\label{def:qm:volatility-models:rv}
If the log price follows $dX_t = \mu_t\,dt + \sigma_t\,dW_t$, the \emph{integrated variance}\index{integrated variance} of day $t$ is $\mathrm{IV}_t = \int_{t-1}^t\sigma_s^2\,ds$.
The \emph{realised variance}\index{realised variance} is the sum of the day's squared intraday returns, $\mathrm{RV}_t = \sum_{j=1}^m r_{t,j}^2$ over $m$ intervals.
\end{definition}

\begin{proposition}[Realised variance converges to integrated variance]\label{prop:qm:volatility-models:rv}
Without microstructure noise, $\mathrm{RV}_t \to \mathrm{IV}_t$ in probability as $m \to \infty$; if $\sigma$ is constant within the day, $\mathrm{RV}_t/\mathrm{IV}_t$ has
mean one and standard deviation $\sqrt{2/m}$.
\end{proposition}

\begin{proof}
$\mathrm{RV}_t$ is the discretised \omterm{def:qm:brownian-motion:qv}{quadratic variation} of $X$, which converges to $\int\sigma^2$ (chapter~3). With constant $\sigma$, each $r_{t,j} \sim \mathcal N(0, \mathrm{IV}_t/m)$
(the drift is negligible at this scale), so $m\,\mathrm{RV}_t/\mathrm{IV}_t$ is $\chi^2_m$, with mean $m$ and variance $2m$.
\end{proof}

A squared daily return is \omterm{def:qm:volatility-models:rv}{realised variance} with $m = 1$: an unbiased but very noisy measure of the day's variance, with a relative error of $\sqrt2 = 141\%$.
Five-minute returns over a 6.5-hour session ($m = 78$) cut it to 16\%. The chapter's simulated market has three volatility components (daily AR coefficients
0.3, 0.9 and 0.99); the measured relative errors are 144\%, 16\% and 7.3\% for $m = 1$, 78 and 390, against the theory's 141\%, 16\% and 7.2\%
(\cref{fig:qm:volatility-models:rv}). In real markets the gain stops at the frequency where bid-ask bounce and discreteness dominate the returns, the
subject of chapter~21.

\begin{omfigure}
\begin{tikzpicture}
\begin{loglogaxis}[width=9cm,height=5cm,xlabel={intraday returns per day, $m$},ylabel={sd of RV/IV $- 1$},
  tick label style={font=\small},label style={font=\small},xmin=0.8,xmax=500,ymin=0.05,ymax=2,grid=major,grid style={gray!25},
  xtick={1,6,13,26,78,390},xticklabels={1,6,13,26,78,390},
  legend style={font=\footnotesize,at={(0.97,0.97)},anchor=north east,fill=white,draw=none},table/col sep=comma]
\addplot[only marks,mark=*,mark size=2pt,omDef] table[x=m,y=error]{figdata/methods/18-volatility-models/rv.csv};
\addplot[omThm,thick,dashed] table[x=m,y=theory]{figdata/methods/18-volatility-models/rv.csv};
\legend{simulated market,$\sqrt{2/m}$}
\end{loglogaxis}
\end{tikzpicture}

\omcaption{Relative error of \omterm{def:qm:volatility-models:rv}{realised variance} as a measure of \omterm{def:qm:volatility-models:rv}{integrated variance}, against the number of intraday returns (1 = daily squared return, 78 =
five minutes over 6.5 hours, 390 = one minute), in a simulated three-component stochastic-volatility market without microstructure noise, and the theory
$\sqrt{2/m}$. Data: the chapter's tutorial, seeded.}\label{fig:qm:volatility-models:rv}
\end{omfigure}

\section{Heterogeneous autoregressive models}

\begin{definition}[HAR model]\label{def:qm:volatility-models:har}
The \emph{HAR model}\index{HAR model} (Corsi, 2009) regresses tomorrow's \omterm{def:qm:volatility-models:rv}{realised variance} on today's, on the average of the last five days and on the average of
the last twenty-two: $\mathrm{RV}_{t+1} = b_0 + b_d\mathrm{RV}_t + b_w\mathrm{RV}_t^{(5)} + b_m\mathrm{RV}_t^{(22)} + e_{t+1}$.
\end{definition}

Three horizons of averaging, fitted by least squares, reproduce the slow hyperbolic-looking decay of volatility autocorrelations with an AR model of order 22
restricted to three parameters: \omterm{def:qm:linear-time-series:longmem}{long memory} approximated by a cascade of short ones. On 2\,000 days of the simulated market the fit is $b_d = 0.35$, $b_w =
0.27$, $b_m = 0.20$ with $R^2 = 0.37$; over the next 1\,000 days its forecasts beat an AR(1) on \omterm{def:qm:volatility-models:rv}{realised variance} (QLIKE $-0.479$ against $-0.473$,
Diebold--Mariano statistic $-4.5$) and are closer to the true \omterm{def:qm:volatility-models:rv}{integrated variance} (mean squared error 0.0073 against 0.0079).

\section{Forecasting and evaluation}

The variance of a day is never observed, only a proxy: a squared return, or a \omterm{def:qm:volatility-models:rv}{realised variance}. A loss function compares a forecast $h$ with the proxy
$\hat\sigma^2$; Patton (2011) showed that the ranking of forecasts is preserved under a noisy but unbiased proxy only for a family of losses that includes the
squared error and QLIKE.

\begin{definition}[QLIKE loss, Mincer--Zarnowitz regression, Diebold--Mariano test]\label{def:qm:volatility-models:eval}
The \emph{QLIKE loss}\index{QLIKE loss} is $\hat\sigma^2/h + \ln h$ (up to terms that do not involve $h$), the negative Gaussian log-likelihood of the proxy's
return: minimised in expectation by $h = \E[\hat\sigma^2]$ and heavier on under-prediction than on over-prediction. The \emph{Mincer--Zarnowitz
regression}\index{Mincer--Zarnowitz regression} (1969) of the proxy on the forecast, $\hat\sigma^2_t = a + bh_t + e_t$, tests unbiasedness ($a = 0$, $b = 1$). The
\emph{Diebold--Mariano test}\index{Diebold--Mariano test} (1995) of equal accuracy divides the mean loss difference of two forecasts by its HAC \omterm{def:qm:estimation:estimator}{standard error}.
\end{definition}

On EUR/USD, the GARCH-$t$ is refitted on 1999--2014 only (persistence 0.9982, \omterm{def:qm:stochastic-differential-equations:ou}{half-life} 392 days) and run forward with fixed parameters; EWMA and the window
need no fitting. Over the 3\,002 test days from 2015 to 2026, with the squared return as proxy, the mean QLIKE is $-0.521$ for GARCH, $-0.495$ for EWMA
and $-0.435$ for the window (\cref{fig:qm:volatility-models:eval}). GARCH beats EWMA with a Diebold--Mariano statistic of $-2.76$ ($p = 0.006$) and the window
with $-4.10$; EWMA beats the window with $-2.48$. The Mincer--Zarnowitz slopes are 0.77, 0.60 and 0.56, all below one (every forecast over-reacts to recent
data), and the $R^2$ are 0.030, 0.027 and 0.010: a squared daily return is so noisy that even a good forecast explains 3\% of it.

\begin{omfigure}
\begin{tikzpicture}
\begin{axis}[width=9cm,height=4.2cm,xbar,bar width=10pt,
  xlabel={mean QLIKE minus the 250-day window's (2015--2026)},
  ytick={0,1,2},yticklabels={GARCH-$t$,{EWMA, 0.94},250-day window},
  tick label style={font=\small},label style={font=\small},
  xmin=-0.1,xmax=0.02,ymin=-0.6,ymax=2.6,grid=major,grid style={gray!25},
  xticklabel style={/pgf/number format/fixed,/pgf/number format/precision=2},scaled x ticks=false,table/col sep=comma]
\addplot[fill=omDef!60,draw=omDef,nodes near coords,nodes near coords style={font=\footnotesize,anchor=west,xshift=2pt,
  /pgf/number format/fixed,/pgf/number format/precision=3}]
  table[x=qlike,y=k]{figdata/methods/18-volatility-models/eval.csv};
\end{axis}
\end{tikzpicture}

\omcaption{Out-of-sample accuracy of three one-day EUR/USD variance forecasts, 2015--2026 (3\,002 days), as mean QLIKE relative to the 250-day window (lower
is better): GARCH-$t$ fitted on 1999--2014 $-0.086$, EWMA $-0.060$. Both differences are significant by the \omterm{def:qm:volatility-models:eval}{Diebold--Mariano test}. Data: ECB euro reference
rates (source: ECB statistics).}\label{fig:qm:volatility-models:eval}
\end{omfigure}

\section{Tutorial: the ghost in the risk report}\label{tut:qm:volatility-models:ghost}

\begin{tutorial}
\textbf{Goal.} Fit GARCH to EUR/USD, find the day a large return leaves the 250-day window, and compare three forecasts out of sample.
\textbf{End state:} \cref{fig:qm:volatility-models:impulse,fig:qm:volatility-models:ghost,fig:qm:volatility-models:eval}, the \omterm{def:qm:stochastic-differential-equations:ou}{half-life} of 445 days and the 7.8\%
drop.
\begin{tutsteps}
\item \textbf{The GARCH recursion and its likelihood} with Gaussian or Student-$t$ innovations.
\omcode{code/firm/volfcst/firm_volfcst.py}{25}{38}{The GJR-GARCH variance recursion and the per-day log-likelihood.}\label{lst:qm:volatility-models:filter}
\item \textbf{Loss functions and the \omterm{def:qm:volatility-models:eval}{Diebold--Mariano test}.}
\omcode{code/firm/volfcst/firm_volfcst.py}{186}{215}{QLIKE, squared error, Mincer--Zarnowitz and Diebold--Mariano.}\label{lst:qm:volatility-models:dm}
\item \textbf{Run} \texttt{fits()}, \texttt{ghost()}, \texttt{evaluate()} and \texttt{har\_vs\_ar1()} in \texttt{qm\_vol.py}, then \texttt{fig\_vol.py}.
\end{tutsteps}
\textbf{What to change next.} Replace the 250-day window by one that downweights returns linearly toward the edge and watch the ghost fade; evaluate the forecasts
at a ten-day horizon, where GARCH's \omterm{def:qm:stochastic-differential-equations:ou}{mean reversion} should matter more.
\end{tutorial}

\section{Build: the volatility forecasting module}\label{bld:qm:volatility-models:volfcst}

\begin{build}
\bfield{Purpose} Every volatility the miniature firm uses for sizing, limits and option quotes is a forecast from this module, evaluated against the others
every month.
\bfield{Interface} \texttt{garch\_filter}, \texttt{garch\_fit(r, dist, gjr)}, \texttt{garch\_forecast(fit, horizon)}; \texttt{ewma(r, lam)}, \texttt{rolling\_var(r,
window)}; \texttt{har\_fit(rv)}, \texttt{har\_forecast(fit, rv)}; \texttt{qlike}, \texttt{mse}, \texttt{mincer\_zarnowitz}, \texttt{diebold\_mariano}.
\bfield{Rules} Forecasts for day $t$ use data up to $t - 1$ only; \omterm{def:qm:estimation:estimator}{standard errors} are the sandwich; persistence and \omterm{def:qm:stochastic-differential-equations:ou}{half-life} are reported with every fit; forecast
comparisons use QLIKE or squared error against an unbiased proxy, with a \omterm{def:qm:volatility-models:eval}{Diebold--Mariano test}.
\bfield{Acceptance tests} \texttt{code/firm/volfcst/tests/}: GARCH and GARCH-$t$ recover simulated parameters; the recursions of GARCH, GJR, EWMA and the window
by hand; the variance forecast's \omterm{def:qm:stochastic-differential-equations:ou}{mean reversion}; QLIKE is minimised by the true variance under a noisy proxy; Mincer--Zarnowitz slope one for the true
variance; the \omterm{def:qm:volatility-models:eval}{Diebold--Mariano test} detects a shifted loss; HAR recovers its coefficients.
\bfield{Stretch} EGARCH; GARCH with \omterm{def:qm:volatility-models:rv}{realised variance} as an extra regressor; multi-day QLIKE evaluation with overlapping horizons.
\end{build}

\begin{omsources}
\item US Bureau of Labor Statistics, \emph{Consumer Price Index -- October 2022}, news release of 10 November 2022; ECB, euro foreign exchange reference rates (concertation around 14:10 CET).
\item R.~F.~Engle, ``Autoregressive conditional heteroscedasticity with estimates of the variance of United Kingdom inflation'', \emph{Econometrica} 50, 1982.
\item T.~Bollerslev, ``Generalized autoregressive conditional heteroskedasticity'', \emph{Journal of Econometrics} 31, 1986.
\item L.~R.~Glosten, R.~Jagannathan and D.~E.~Runkle, ``On the relation between the expected value and the volatility of the nominal excess return on stocks'',
\emph{Journal of Finance} 48, 1993.
\item J.P.~Morgan/Reuters, \emph{RiskMetrics --- Technical Document}, 4th ed., 1996.
\item T.~G.~Andersen, T.~Bollerslev, F.~X.~Diebold and P.~Labys, ``Modeling and forecasting realized volatility'', \emph{Econometrica} 71, 2003;
O.~E.~Barndorff-Nielsen and N.~Shephard, \emph{Journal of the Royal Statistical Society B} 64, 2002.
\item F.~Corsi, ``A simple approximate long-memory model of realized volatility'', \emph{Journal of Financial Econometrics} 7, 2009.
\item A.~J.~Patton, ``Volatility forecast comparison using imperfect volatility proxies'', \emph{Journal of Econometrics} 160, 2011.
\item F.~X.~Diebold and R.~S.~Mariano, ``Comparing predictive accuracy'', \emph{Journal of Business and Economic Statistics} 13, 1995; J.~Mincer and
V.~Zarnowitz, ``The evaluation of economic forecasts'', NBER, 1969.
\item European Central Bank, euro foreign exchange reference rate against the dollar (ECB Data Portal, series \texttt{EXR.D.USD.EUR.SP00.A}), accessed 24
September 2026.
\end{omsources}

\section{Exercises}

\begin{exercise}[$\star$]\label{exo:qm:volatility-models:1}
A GARCH(1,1) has $\omega = 0.02$, $\alpha = 0.05$ and $\beta = 0.93$ (daily, returns in percent). What are its unconditional daily and annual volatilities and
the \omterm{def:qm:stochastic-differential-equations:ou}{half-life} of a variance shock?
\end{exercise}

\begin{exercise}[$\star$]\label{exo:qm:volatility-models:2}
With $\lambda = 0.94$, after how many days does a squared return's weight in the EWMA fall by half? What share of the total weight do the last 20 days carry?
\end{exercise}

\begin{exercise}[$\star$]\label{exo:qm:volatility-models:3}
\omterm{def:qm:volatility-models:rv}{Realised variance} from 5-minute returns over 24 hours: what is its relative \omterm{def:qm:estimation:estimator}{standard error} as a measure of \omterm{def:qm:volatility-models:rv}{integrated variance}, without noise?
\end{exercise}

\begin{exercise}[$\star\star$]\label{exo:qm:volatility-models:4}
In the same GARCH(1,1), today's conditional variance is four times the unconditional. What is the forecast of the average variance over the next 10 days?
\end{exercise}

\begin{exercise}[$\star\star$]\label{exo:qm:volatility-models:5}
Show that QLIKE, $\hat\sigma^2/h + \ln h$, is minimised in expectation at $h = \E[\hat\sigma^2]$, and compare the losses of forecasts $0.5\sigma^2$ and $1.5\sigma^2$.
\end{exercise}

\begin{exercise}[$\star\star$]\label{exo:qm:volatility-models:6}
A 250-day window contains one return of 3.5\% among days of 0.5\%. By what percentage does the window's volatility fall when it leaves?
\end{exercise}

\begin{exercise}[$\star\star\star$]\label{exo:qm:volatility-models:7}
\emph{Coding.} Fit GARCH(1,1) to the EUR/USD returns with the normal likelihood and report both the Hessian and the sandwich \omterm{def:qm:estimation:estimator}{standard errors} of
$\alpha$ and $\beta$. Which should be reported, and why do they differ?
\end{exercise}

\begin{exercise}[$\star\star\star$]\label{exo:qm:volatility-models:8}
\emph{Find the flaw.} ``We compared our volatility forecasts by their mean absolute error against the absolute daily return; the EWMA with $\lambda = 0.97$ wins.''
\end{exercise}

\section{Problem: The Ghost in the Risk Report}

\begin{problem}[{Weekend problem --- a volatility shock and the calendar}]\label{pb:qm:volatility-models:1}
The data are the ECB's daily euro reference rates against the dollar, 4 January 1999 to 23 September 2026, as log returns in percent (7\,098 days).

\medskip\noindent\textbf{Part I --- The fit.}
\begin{enumerate}
\item What GARCH(1,1) parameters does the normal likelihood give, and their sandwich \omterm{def:qm:estimation:estimator}{standard errors}?
\item And the Student-$t$ likelihood? How much higher is its log-likelihood?
\item What are the persistence and the \omterm{def:qm:stochastic-differential-equations:ou}{half-life} of a variance shock?
\item What is the unconditional daily volatility, and how does it compare with the sample standard deviation?
\item Is there a \omterm{def:qm:volatility-models:gjr}{leverage effect}?
\end{enumerate}

\medskip\noindent\textbf{Part II --- The ghost.}
\begin{enumerate}[resume]
\item Which return causes the largest one-day relative fall of the 250-day volatility, and when?
\item By how much does the window's volatility fall, and what did the market do that day?
\item What did the GARCH and EWMA volatilities do on the same day?
\item How long does a 3.5\% return stay in each forecast?
\item What would a report based on the window have said about risk that day?
\end{enumerate}

\medskip\noindent\textbf{Part III --- Which forecast?}
\begin{enumerate}[resume]
\item What are the mean QLIKE losses of the three forecasts from 2015 to 2026?
\item Are the differences significant?
\item What do the Mincer--Zarnowitz slopes say?
\item Why is the $R^2$ of every forecast so low?
\item What would \omterm{def:qm:volatility-models:rv}{realised variance} from intraday data change in the evaluation?
\end{enumerate}

\medskip\noindent\textbf{Part IV --- Judgement.}
\begin{enumerate}[resume]
\item Which volatility should the risk report use?
\item Is a \omterm{def:qm:stochastic-differential-equations:ou}{half-life} of more than a year believable?
\item What would you tell the risk committee about the day of the ghost?
\item State the \emph{named result}: the \omterm{def:qm:stochastic-differential-equations:ou}{half-life} of a EUR/USD volatility shock and the size of the ghost.
\item In one sentence: what is wrong with an equally weighted window?
\end{enumerate}
\end{problem}

\section{Interview questions}

\begin{interviewq}[$\star$ \iqroles{risk, trader}]\label{iq:qm:volatility-models:1}
Why can an equally weighted historical volatility drop on a volatile day?
\end{interviewq}

\begin{interviewq}[$\star\star$ \iqroles{researcher, risk}]\label{iq:qm:volatility-models:2}
Write down GARCH(1,1). What is its unconditional variance, and what does $\alpha + \beta$ mean?
\end{interviewq}

\begin{interviewq}[$\star\star$ \iqroles{researcher}]\label{iq:qm:volatility-models:3}
How is EWMA related to GARCH, and what does it get wrong about long-horizon forecasts?
\end{interviewq}

\begin{interviewq}[$\star\star$ \iqroles{researcher, trader}]\label{iq:qm:volatility-models:4}
How do you evaluate a volatility forecast when the true volatility is never observed?
\end{interviewq}

\begin{interviewq}[$\star\star$ \iqroles{researcher, trader}]\label{iq:qm:volatility-models:5}
What is \omterm{def:qm:volatility-models:rv}{realised variance}, and why not sample it every second?
\end{interviewq}

\begin{interviewq}[$\star\star\star$ \iqroles{researcher}]\label{iq:qm:volatility-models:6}
Why does the \omterm{def:qm:volatility-models:har}{HAR model} forecast volatility well with only three regressors?
\end{interviewq}
